\section{Structure theory}

\begin{definition}
A \emph{torus} $\at \subseteq \g$ is an abelian subalgebra such that for all
$t \in \at$ the adjoint $\ad(t)\colon \g \to \g$ is a diagonalisable linear
map. A \emph{maximal torus} is a torus that is not contained in any bigger
torus.
\end{definition}

\begin{example}
Let $T = (S^1)^r \subseteq G$ where $G$ is a compact Lie group (or
$T = (\C^\times)^r \subseteq G$ where $G$ is a reductive algebraic group). Then
$\at = \Lie(T) \subseteq \g = \Lie(G)$ is a torus, maximal if $T$ is. (This
comes from knowing that representations of $T$ are diagonalisable by an
averaging argument, but we don't really know this.)
\end{example}

\begin{exercise}\leavevmode
\begin{enumerate}
  \item For $\g = \sl_n$ or $\g = \gl_n$, the diagonal matrices (of trace $0$
    if inside $\sl_n$) form a maximal torus $\at$.
  \item $\bigl(\begin{smallmatrix} 0 & * \\ 0 & 0 \end{smallmatrix}\bigr)$ is
    \emph{not} a torus. (We should know this is not diagonalisable!)
\end{enumerate}
\end{exercise}

For a vector space $V$, let $t_1, \dots, t_r\colon V \to V$ be pairwise
commuting diagonalisable linear maps; let
$\lambda = (\lambda_1, \dots, \lambda_r) \in \C^r$. Set
$V_\lambda = \setc{v \in V}{t_iv = \lambda_iv \text{ for all } i}$, the
simultaneous eigenspace.

\begin{lemma}
$V = \bigoplus_{\lambda \in \C^r} V_\lambda$.
\end{lemma}

\begin{proof}
Induct on $r$. For $r = 1$ this follows from the requirement that $t_1$ be
diagonalisable. For $r > 1$, look at $t_1, \dots, t_{r-1}$ and decompose
\[
  V = \bigoplus V_{(\lambda_1, \dots, \lambda_{r-1})}.
\]
Since $t_r$ commutes with $t_1, \dots, t_{r-1}$, it preserves the
decomposition. Now decompose each $V_{(\lambda_1, \dots, \lambda_{r-1})}$ into
eigenspaces for $t_r$ (the restriction of a diagonalisable map to an invariant
subspace is diagonalisable).
\end{proof}

Let $\at$ be the $r$-dimensional abelian Lie algebra with basis
$(t_1, \dots, t_r)$. The lemma asserts that $V$ is a semisimple (i.e.\
completely reducible) representation of $\at$. $V = \bigoplus V_\lambda$ is the
decomposition into isotypic representations. Namely, letting $\C_\lambda$ be
the $1$-dimensional representation wherein $t_i(w) = \lambda_iw$, we have
$\lambda \neq \mu \implies \C_\lambda \not\cong \C_\mu$, and $V_\lambda$ is
the direct sum of $\dim V_\lambda$ copies of $\C_\lambda$.

\begin{exercise}
Show that every irreducible representation of $\at$ is one-dimensional.
\end{exercise}

Another way of saying this is as follows. $\lambda$ is a linear map
$\at \to \C$, i.e.\ an element of the dual space $\lambda \in \at^*$ (sending
$t_i \mapsto \lambda_i$). Therefore, one-dimensional representations of $\at$
(which are all the irreducible representations of $\at$) correspond to
elements of $\at^* = \Hom_{\text{vector spaces}}(\at, \C)$. The decomposition
\[
  V = \bigoplus_{\lambda \in \at^*} V_\lambda, \qquad
  V_\lambda = \setc{v \in V}{t(v) = \lambda(t)v \text{ for all } t \in \at}
\]
is called the \emph{weight space decomposition} of $V$.

Now, let $\g$ be a Lie algebra and $\at$ a maximal torus. The weight space
decomposition of $\g$ is
\[
  \g = \g_0 \oplus \bigoplus_{\substack{\lambda \in \at^*\\ \lambda \neq 0}} \g_\lambda,
\]
where $\g_0 = \setc{x \in \g}{[t, x] = 0 \text{ for all } t \in \at}$ and
$\g_\lambda = \setc{x \in \g}{[t, x] = \lambda(t)x \text{ for all } t \in \at}$.

\begin{definition}
$R = \setc{\lambda \in \at^*}{\g_\lambda \neq 0,\ \lambda \neq 0}$ is the set
of \emph{roots} of $\g$.
\end{definition}

We will now compute the root space decomposition for $\sl_n$.

\begin{example}
$\g = \sl_n$, $\at = $ diagonal matrices in $\sl_n$. Let
\[
  t = \begin{pmatrix} t_1 & & 0 \\ & \ddots & \\ 0 & & t_n \end{pmatrix}
\]
and let $E_{ij}$ be the matrix with $1$ in the $(i, j)$th position and $0$
elsewhere. Then $[t, E_{ij}] = (t_i - t_j)E_{ij}$.

Let $\eps_i \in \at^*$, $\eps_i(t) = t_i$. Then $\eps_1, \dots, \eps_n$ span
$\at^*$, but $\eps_1 + \dots + \eps_n = 0$.

\begin{remark}
$\at$ is a subalgebra of $\C^n$, so $\at^*$ should be a quotient of
$(\C^n)^*$, and here is the relation by which it is a quotient.
\end{remark}

Therefore, we can rewrite the above as
$[t, E_{ij}] = (\eps_i - \eps_j)(t) \cdot E_{ij}$. Therefore, the roots are
$R = \setc{\eps_i - \eps_j}{i \neq j}$ and $\g_0 = \at$. (Note that we've shown
$\at$ is a maximal torus in the process.) The root space
$\g_{\eps_i - \eps_j} = \C E_{ij}$ is one-dimensional.

That is, the root space decomposition of $\sl_n$ is
\[
  \sl_n = \at \oplus \bigoplus_{\eps_i - \eps_j \in R} \g_{\eps_i - \eps_j}.
\]
\end{example}

\begin{exercise}
``This is (a) an essential exercise, and (b) guaranteed to be on the exam. The
first should get you to do it, the second is irrelevant.''

Compute the root space decomposition for $\sl_n$ (as we just did),
$\so_{2n}$, $\so_{2n+1}$, $\spl_{2n}$, where $\at$ is the diagonal matrices
in $\g$, and we take the conjugate $\so_n$ defined by
\[
  \so_n = \setc{A \in \gl_n}{JA + A^TJ = 0}, \qquad
  J = \begin{pmatrix} 0 & & 1 \\ & \iddots & \\ 1 & & 0 \end{pmatrix}
\]
(this is conjugate to the usual $\so_n$ since we've defined a nondegenerate
inner product; in this version, the diagonal matrices are a maximal torus).
The symplectic matrices are given by the same embedding as before.

In particular, show that $\at$ is a maximal torus and the root spaces are
one-dimensional.
\end{exercise}

\lecture{9} The Lie algebras we were working with last time are called the
``classical'' Lie algebras (they are Lie algebras of automorphisms of vector
spaces preserving a form). The rest of the course will be concerned with
explaining how their roots control their representations.

\begin{proposition}
$\sl_n(\C)$ is simple.
\end{proposition}

\begin{proof}
Recall
\[
  \sl_n(\C) = \at \oplus \bigoplus_{\alpha \in R} \g_\alpha,
\]
where $R = \setc{\eps_i - \eps_j}{i \neq j}$ and $\g_{\eps_i - \eps_j} = \C E_{ij}$.

Suppose $\fr \subseteq \sl_n$ is a nonzero ideal. Choose $r \in \fr$,
$r \neq 0$, such that when we write $r = t + \sum_{\alpha \in R} e_\alpha$
with $t \in \at$ and $e_\alpha \in \g_\alpha$, the number of nonzero terms is
minimal.

First, suppose $t \neq 0$. Choose $t_0 \in \at$ such that $\alpha(t_0) \neq 0$
for all $\alpha \in R$ (i.e., $t_0$ has distinct eigenvalues). Consider
$[t_0, r] = \sum_{\alpha \in R} \alpha(t_0)e_\alpha$. Note that this lies in
$\fr$, and if it is nonzero, it has fewer nonzero terms than $r$, a
contradiction. Thus, $[t_0, r] = 0$, i.e.\ $r = t \in \at$.

Since $t \neq 0$, there exists $\alpha = \eps_i - \eps_j \in R$ with
$\alpha(t) \neq 0$ (i.e.\ $t$ can't have all its eigenvalues equal, since it
has trace $0$). Then
\[
  [t, E_{ij}] = \alpha(t)E_{ij} \in \fr \implies E_{ij} \in \fr.
\]
Now,
\[
  [E_{ij}, E_{jk}] = E_{ik} \text{ for } i \neq k, \qquad
  [E_{si}, E_{ij}] = E_{sj} \text{ for } s \neq j,
\]
and therefore $E_{ab} \in \fr$ for all $a \neq b$. Finally,
\[
  [E_{i,i+1}, E_{i+1,i}] = E_{ii} - E_{i+1,i+1} \in \fr,
\]
so in fact the entire basis for $\sl_n$ lies in $\fr$ and $\fr = \sl_n$.

\begin{remark}
This is actually a combinatorial statement about root spaces; we'll see more
about this later.
\end{remark}

That leaves us with the case $t = 0$. If $r = cE_{ij}$ has only one nonzero
term, we are done exactly as above; therefore,
\[
  r = e_\alpha + e_\beta + \sum_{\gamma \in R - \{\alpha, \beta\}} e_\gamma,
  \qquad \alpha \neq \beta,\ e_\alpha, e_\beta \neq 0.
\]
Choose $t_0 \in \at$ such that $\alpha(t_0) \neq \beta(t_0)$; then
$[t_0, r] - \beta(t_0)r$ is nonzero and has fewer nonzero terms than $r$, a
contradiction.
\end{proof}

\begin{proposition}
Let $\g$ be a semisimple Lie algebra. Then maximal tori exist (i.e.\ are
nonzero). Moreover,
$\g_0 = \setc{x \in \g}{[\at, x] = 0} = \at$. (Recall that $\at$ is not defined
to be the maximal abelian subalgebra, but as the maximal abelian subalgebra
whose adjoint acts semisimply.)
\end{proposition}

We won't prove this; the proof involves introducing the Cartan subalgebra and
showing that it's the same as a maximal torus. As a result, we will sometimes
be calling $\at$ the Cartan subalgebra.

This means that the root space decomposition of a semisimple Lie algebra is
\[
  \g = \at \oplus \bigoplus_{\alpha \in R} \g_\alpha,
\]
as we have, of course, already seen for the classical Lie algebras.

\begin{theorem}[Structure theorem for semisimple Lie algebras, part I]
Let $\g$ be a semisimple Lie algebra over $\C$, and let $\at \subseteq \g$ be a
maximal torus. Write $\g = \at \oplus \bigoplus_{\alpha \in R} \g_\alpha$. Then
\begin{enumerate}
  \item $\C R = \at^*$, i.e.\ the roots span $\at^*$.
  \item $\dim \g_\alpha = 1$.
  \item If $\alpha, \beta \in R$ and $\alpha + \beta \in R$ then
    $[\g_\alpha, \g_\beta] = \g_{\alpha+\beta}$. If $\alpha + \beta \notin R$
    and $\alpha \neq -\beta$ then $[\g_\alpha, \g_\beta] = 0$.
  \item $[\g_\alpha, \g_{-\alpha}] \subseteq \at$ is one-dimensional, and
    $\g_\alpha \oplus [\g_\alpha, \g_{-\alpha}] \oplus \g_{-\alpha}$ is a Lie
    subalgebra of $\g$ isomorphic to $\sl_2$. (In particular,
    $\alpha \in R \implies -\alpha \in R$.)
\end{enumerate}
\end{theorem}

\begin{exercise}
Check this for the classical Lie algebras.
\end{exercise}

\begin{proof}\leavevmode
(1) Suppose not; then there exists a nonzero $t \in \at$ such that
$\alpha(t) = 0$ for all $\alpha \in R$. But then $[t, \g_\alpha] = 0$; as
$[t, \at] = 0$ by definition, we see that $t$ is central in $\g$. However, $\g$
is semisimple, so it has no abelian ideals, in particular no centre.

We will now show a sequence of statements.
\begin{alist}
  \item $[\g_\lambda, \g_\mu] \subseteq \g_{\lambda+\mu}$ if
    $\lambda, \mu \in \at^*$. Indeed, for $x \in \g_\lambda$ and
    $y \in \g_\mu$ we have
    \[
      [t, [x, y]] = [[t, x], y] + [x, [t, y]] = (\lambda(t) + \mu(t))[x, y].
    \]
    This shows that $[\g_\alpha, \g_\beta] \subseteq \g_{\alpha+\beta}$ if
    $\alpha + \beta \in R$, is $0$ if $\alpha + \beta \notin R$ and
    $\alpha \neq -\beta$, and $[\g_\alpha, \g_{-\alpha}] \subseteq \at$.
  \item $(\g_\lambda, \g_\mu)_{\ad} = 0$ if $\lambda \neq -\mu$;
    $(\,,\,)_{\ad}|_{\g_\lambda \oplus \g_{-\lambda}}$ is nondegenerate.

    \begin{proof}
    Let $x \in \g_\lambda$ and $y \in \g_\mu$. Recall
    $(x, y)_{\ad} = \tr_\g(\ad x \cdot \ad y)$. To show that it's zero, we
    show that $\ad x \cdot \ad y$ is nilpotent. Indeed,
    $(\ad x \cdot \ad y)^N\g_\alpha \subseteq \g_{\alpha + N(\lambda+\mu)}$,
    so if $\lambda \neq -\mu$ then for $N \gg 0$ this is zero.

    On the other hand, $(\,,\,)_{\ad}$ is nondegenerate and
    $\g = \bigoplus_{\lambda \in \at^*} \g_\lambda$, so the restriction
    $(\,,\,)_{\ad}|_{\g_\lambda \oplus \g_{-\lambda}}$ must be nondegenerate.
    \end{proof}
  \item In particular, $(\,,\,)_{\ad}|_\at$ is nondegenerate, as
    $\at = \g_0$.

    \begin{remark}
    $(\,,\,)_{\ad}|_\at \neq (\,,\,)_{\ad\,\at}$ (which is zero since $\at$ is
    abelian).
    \end{remark}

    Therefore, the Killing form defines an isomorphism
    $\nu\colon \at \to \at^*$ with $\nu(t)(t') = (t, t')_{\ad}$. It also
    defines an induced inner product (well, symmetric bilinear form to be
    precise; we don't know it's positive) on $\at^*$, via
    $(\nu(t), \nu(t')) = (t, t')_{\ad}$.
  \item $\alpha \in R \implies -\alpha \in R$, since $(\,,\,)_{\ad}$ is
    nondegenerate on $\g_\alpha \oplus \g_{-\alpha}$, but
    $(\g_\alpha, \g_\alpha)_{\ad} = 0$ if $\alpha \neq 0$. In particular,
    $\g_{-\alpha} \cong \g_\alpha^*$ via the Killing form.
  \item $x \in \g_\alpha$, $y \in \g_{-\alpha} \implies
    [x, y] = (x, y)_{\ad}\, \nu^{-1}(\alpha)$.

    \begin{proof}
    For $t \in \at$,
    \[
      (t, [x, y])_{\ad} = ([t, x], y)_{\ad} = \alpha(t)(x, y)_{\ad}
        = \bigl(t, (x, y)_{\ad}\,\nu^{-1}(\alpha)\bigr)_{\ad},
    \]
    which is exactly what we want, since $(\,,\,)_{\ad}|_\at$ is
    nondegenerate.
    \end{proof}
  \item Let $e_\alpha \in \g_\alpha$ be nonzero, and pick
    $e_{-\alpha} \in \g_{-\alpha}$ such that $(e_\alpha, e_{-\alpha})_{\ad} \neq 0$.
    Then $[e_\alpha, e_{-\alpha}] = (e_\alpha, e_{-\alpha})_{\ad}\,\nu^{-1}(\alpha)$.
    To show that $e_\alpha$, $e_{-\alpha}$, and their bracket give a copy of
    $\sl_2$, we need to compute
    \[
      [\nu^{-1}(\alpha), e_\alpha] = \alpha(\nu^{-1}(\alpha))e_\alpha = (\alpha, \alpha)e_\alpha.
    \]
    We will be done if we show that $(\alpha, \alpha) \neq 0$ (then we get a
    copy of $\sl_2$ after renormalising).

    \begin{proposition}
    $(\alpha, \alpha) \neq 0$ for all $\alpha \in R$.
    \end{proposition}

    \begin{proof}
    Next time.
    \end{proof}
\end{alist}
\let\qed\relax
\end{proof}

\lecture{10} We now prove the last proposition, $(\alpha, \alpha) \neq 0$ for
all $\alpha \in R$.

\begin{proof}
Suppose otherwise. Let
$\fm_\alpha = \langle e_\alpha, e_{-\alpha}, \nu^{-1}(\alpha)\rangle$. If
$(\alpha, \alpha) = 0$, i.e.\ if $[\nu^{-1}(\alpha), e_\alpha] = 0$, then
$[\fm_\alpha, \fm_\alpha] = \C\nu^{-1}(\alpha)$ and $\fm_\alpha$ is solvable.
However, by Lie's theorem this implies that $\ad[\fm_\alpha, \fm_\alpha]$ acts
by nilpotents on $\g$, i.e.\ $\ad\nu^{-1}(\alpha)$ is nilpotent. Since
$\nu^{-1}(\alpha) \in \at$, it is also diagonalisable, and the only
diagonalisable nilpotent element is $0$. However,
$\alpha \in R \implies \alpha \neq 0$.
\end{proof}

Now, define
\[
  h_\alpha = \frac{2}{(\alpha, \alpha)}\nu^{-1}(\alpha)
\]
and rescale $e_{-\alpha}$ such that $(e_\alpha, e_{-\alpha})_{\ad} = 2/(\alpha, \alpha)$.

\begin{exercise}
The map $(e_\alpha, h_\alpha, e_{-\alpha}) \mapsto (e, h, f)$ gives an
isomorphism $\fm_\alpha \cong \sl_2$.
\end{exercise}

\begin{remark}
In $\sl_n$, the root spaces are spanned by the $E_{ij}$, so we are saying that
the subalgebra of matrices of the form
\[
  \begin{pmatrix} a & * \\ * & -a \end{pmatrix}
\]
(entries in rows and columns $i$ and $j$ only, zero elsewhere) is a copy of
$\sl_2$,
which is obvious. The cool thing is that something like this happens for all
the semisimple Lie algebras, and that this lets us know just about everything
about them.
\end{remark}

We will now use our knowledge of the representation theory of $\sl_2$. We show

\begin{claim}
$\dim \g_{-\alpha} = 1$ for all $\alpha \in R$.
\end{claim}

\begin{proof}
Pick a copy of $\fm_\alpha = \langle e_\alpha, h_\alpha, e_{-\alpha}\rangle \cong \sl_2$.
Suppose $\dim \g_{-\alpha} > 1$; then the map $\g_{-\alpha} \to \C\nu^{-1}(\alpha)$,
$x \mapsto [e_\alpha, x]$ has a kernel. That is, there is a nonzero
$v \in \g_{-\alpha}$ such that $\ad e_\alpha \cdot v = 0$. Then $v$ is a
highest-weight vector (it comes from a root space $\g_{-\alpha}$, so it's an
eigenvector of $\ad h_\alpha$), and its weight is determined by
\[
  \ad h_\alpha \cdot v = -\alpha(h_\alpha)v = -2v.
\]
That is, $v$ is a highest-weight vector of weight $-2$. However, we know that
in finite-dimensional representations of $\sl_2$ (which $\g$ is, since
$\fm_\alpha$ acts on it) the highest weights are nonnegative integers. This
contradiction shows the claim.
\end{proof}

Before we finish proving the theorem we had (we still need to show
$[\g_\alpha, \g_\beta] = \g_{\alpha+\beta}$ when
$\alpha, \beta, \alpha + \beta \in R$), we show further combinatorial
properties of the roots.

\begin{theorem}[Structure theorem, continued]\leavevmode
\begin{enumerate}
  \item $\dfrac{2(\alpha, \beta)}{(\alpha, \alpha)} \in \Z$ for all
    $\alpha, \beta \in R$.
  \item If $\alpha \in R$ and $k\alpha \in R$, then $k = \pm 1$.
  \item Let $\alpha, \beta \in R$ with $\beta \neq \pm\alpha$. Then
    \[
      \bigoplus_{k \in \Z} \g_{\beta + k\alpha}
    \]
    is an irreducible module for
    $(\sl_2)_\alpha = \langle e_\alpha, h_\alpha, e_{-\alpha}\rangle$. In
    particular,
    \[
      \setc{\beta + k\alpha}{\beta + k\alpha \in R,\ k \in \Z}
    \]
    is of the form
    $\beta - p\alpha, \beta - (p - 1)\alpha, \dots, \beta, \dots,
    \beta + (q - 1)\alpha, \beta + q\alpha$, where
    $p - q = \frac{2(\alpha, \beta)}{(\alpha, \alpha)}$. We call this set the
    ``$\alpha$-string through $\beta$''.
\end{enumerate}
\end{theorem}

\begin{proof}
(1): Let $q = \max\setc{k}{\beta + k\alpha \in R}$ and let
$v \in \g_{\beta + q\alpha}$ be nonzero. Then
$\ad e_\alpha \cdot v \in \g_{\beta + (q+1)\alpha} = 0$, and
\[
  \ad h_\alpha \cdot v = (\beta + q\alpha)(h_\alpha) \cdot v
    = \Bigl(\frac{2(\beta, \alpha)}{(\alpha, \alpha)} + 2q\Bigr)v.
\]
That is, $v$ is a highest-weight vector for $(\sl_2)_\alpha$ with weight
$2(\beta, \alpha)/(\alpha, \alpha) + 2q$. Since the weight in a
finite-dimensional representation is an integer, we see that
$2(\beta, \alpha)/(\alpha, \alpha) \in \Z$.

(3): The structure of $\sl_2$-modules tells us that
$(\ad e_{-\alpha})^rv \neq 0$ for
$0 \leq r \leq \frac{2(\alpha, \beta)}{(\alpha, \alpha)} + 2q = N$, and that
$(\ad e_{-\alpha})^{N+1}v = 0$. Therefore,
$\setc{\beta + q\alpha - k\alpha}{0 \leq k \leq N}$ are all roots (or possibly
equal to zero; in any case, they have nonzero root space). This certainly is
an irreducible $(\sl_2)_\alpha$-module; we now need to show that there are no
other roots in the $\alpha$-string through $\beta$. We do this by repeating the
same construction from the bottom up.

Let $p = \max\setc{k}{\beta - k\alpha \in R}$ and let
$w \in \g_{\beta - p\alpha}$ be nonzero. Then $\ad e_{-\alpha} \cdot w = 0$ and
$\ad h_\alpha \cdot w = (2(\beta, \alpha)/(\alpha, \alpha) - 2p)w$, so $w$ is a
lowest-weight vector of weight $2(\beta, \alpha)/(\alpha, \alpha) - 2p$. By
applying $\ad e_\alpha$ repeatedly, we get an irreducible $\sl_2$-module with
roots $\beta - p\alpha, \beta - (p - 1)\alpha, \dots,
\beta + (p - 2(\alpha, \beta)/(\alpha, \alpha))\alpha$.

Now by construction we have $p - 2(\alpha, \beta)/(\alpha, \alpha) \leq q$ and
also $q + 2(\alpha, \beta)/(\alpha, \alpha) \leq p$, which means that in fact
$p - q = 2(\alpha, \beta)/(\alpha, \alpha)$ and the two submodules we get
coincide.

(2): By (1) we have $2(\alpha, k\alpha)/(k\alpha, k\alpha) = 2/k \in \Z$ and
$2(k\alpha, \alpha)/(\alpha, \alpha) = 2k \in \Z$. Thus, it suffices to show
that $\alpha \in R \implies 2\alpha \notin R$.

But indeed, suppose $2\alpha \in R$ and let $v \in \g_{-2\alpha}$ be nonzero.
Then $[e_\alpha, v] \in \g_{-\alpha}$ has the property that
\[
  ([e_\alpha, v], e_\alpha) = -(v, [e_\alpha, e_\alpha]) = 0,
\]
and since $\g_\alpha$ is one-dimensional and spanned by $e_\alpha$, we find
$([e_\alpha, v], \g_\alpha) = 0$. Since
$(\,,\,)|_{\g_\alpha \oplus \g_{-\alpha}}$ is nondegenerate, this means
$[e_\alpha, v] = 0$, i.e.\ $v \in \g_{-2\alpha}$ is a highest-weight vector
for $(\sl_2)_\alpha$ with highest weight $-4$. This doesn't happen in
finite-dimensional representations, so $\g_{-2\alpha} = 0$ and $2\alpha$ is not
a root.

(In fact, we could derive this from (3), since we know that there is a
$3$-dimensional $\sl_2$-module running through $0$.)
\end{proof}

We finally show $[\g_\alpha, \g_\beta] = \g_{\alpha+\beta}$ when
$\alpha, \beta, \alpha + \beta \in R$. Indeed, we have just shown that
$\bigoplus_k \g_{\beta + k\alpha}$ is an irreducible representation of
$\fm_\alpha$, and $\ad e_\alpha\colon \g_{\beta + k\alpha} \to
\g_{\beta + (k+1)\alpha}$ is an isomorphism for $-p \leq k < q$. Since
$q \geq 1$, we see that $\ad e_\alpha \cdot \g_\beta = \g_{\alpha+\beta}$.

\begin{remark}
The combinatorial statement of (3) is rather clunky. We now ``unclunk'' it.

For $\alpha \in \at^*$, define $s_\alpha\colon \at^* \to \at^*$ by
$s_\alpha(v) = v - \frac{2(\alpha, v)}{(\alpha, \alpha)}\alpha$ (a reflection
of $v$ in the hyperplane orthogonal to $\alpha$).
\end{remark}

\begin{claim}
(3) implies (and in fact, is equivalent to) $s_\alpha(\beta) \in R$ whenever
$\alpha, \beta \in R$.
\end{claim}

\begin{proof}
For $\beta = \pm\alpha$ this is clear, so let $\beta \neq \pm\alpha$. Let
$r = 2(\alpha, \beta)/(\alpha, \alpha)$. If $r \geq 0$, then
$p = q + r \geq r$; if $r \leq 0$ then $q = p - r \geq -r$. Either way,
$\beta - r\alpha$ is in the $\alpha$-string through $\beta$.
\end{proof}

Over the next several lectures, we will classify the root systems as
combinatorial objects. Before we do that, we state some further properties of
the roots:

\begin{proposition}\leavevmode
\begin{enumerate}
  \item $\alpha, \beta \in R \implies (\alpha, \beta) \in \Q$.
  \item If we pick a basis $\beta_1, \dots, \beta_l$ of $\at^*$, with
    $\beta_i \in R$, then any $\beta \in R$ expands as $\sum q_i\beta_i$ with
    $q_i \in \Q$, i.e.\ $\dim_\Q \Q R = \dim_\C \at$.
  \item $(\,,\,)$ is positive definite on $\Q R$.
\end{enumerate}
\end{proposition}

\lecture{11}

\begin{proof}[Proof of the proposition at the end of last lecture]\leavevmode
\begin{enumerate}
  \item It suffices to show $(\beta, \beta) \in \Q$ for all $\beta \in R$.
    Let $h, h' \in \at$; then
    \[
      (h, h')_{\ad} = \tr_\g(\ad h \ad h') = \sum_{\alpha \in R} \alpha(h)\alpha(h'),
    \]
    since $\g = \at + \bigoplus_{\alpha \in R} \g_\alpha$ by the structure
    theorem. Therefore, for $\lambda, \mu \in \at^*$,
    \[
      (\lambda, \mu) = (\nu^{-1}(\lambda), \nu^{-1}(\mu))_{\ad}
        = \sum_{\alpha \in R} \alpha(\nu^{-1}\lambda)\alpha(\nu^{-1}\mu)
        = \sum_{\alpha \in R} (\lambda, \alpha)(\mu, \alpha).
    \]
    In particular, $(\beta, \beta) = \sum_{\alpha \in R} (\alpha, \beta)^2$.
    Multiplying by $4/(\beta, \beta)^2$, we get
    \[
      \frac{4}{(\beta, \beta)} = \sum_{\alpha \in R}
        \Bigl(\frac{2(\alpha, \beta)}{(\beta, \beta)}\Bigr)^2 \in \Z,
    \]
    and therefore $(\beta, \beta) \in \Q$.
  \item Let $B$ be the Gram matrix of the basis, $B_{ij} = (\beta_i, \beta_j)$.

    \begin{exercise}
    $(\,,\,)$ is a nondegenerate symmetric bilinear form $\implies \det B \neq 0$.
    \end{exercise}

    Let $\beta = \sum c_i\beta_i \in R$. Then
    $(\beta, \beta_i) = \sum_j c_j(\beta_j, \beta_i)$. Therefore, we can recover
    the vector of $c_i$ as
    \[
      (c_1, \dots, c_l) = \bigl((\beta, \beta_1), \dots, (\beta, \beta_l)\bigr)B^{-1}.
    \]
    Since the inner products of roots are all rational, we see that
    $c_i \in \Q$ as well.
  \item Let $\lambda \in \Q R$. Then $\lambda = \sum c_i\beta_i$ with
    $c_i \in \Q$, and $(\lambda, \alpha) \in \Q$ for all $\alpha \in R$.
    Moreover, $(\lambda, \lambda) = \sum_{\alpha \in R} (\lambda, \alpha)^2 \geq 0$,
    and if $(\lambda, \lambda) = 0$ then $(\lambda, \alpha) = 0$ for all
    $\alpha \in R$. Since $R$ spans $\at^*$ and $(\,,\,)$ is nondegenerate,
    this implies that $\lambda = 0$. \qedhere
\end{enumerate}
\end{proof}
